%! AP Statistics — Unit 9, Lesson 9.4 — Guided Notes (Answer Key)
\documentclass[10pt]{article}
\usepackage{apstats-article}
\usepackage{apstats-key}

\begin{document}

\pageheader{Unit 9, Lesson 9.4}{Guided Notes --- Answer Key}
\namedateperiod

\begin{objectivebox}
By the end of this lesson, I will be able to:
\begin{itemize}
    \item[$\square$] \writeline
    \item[$\square$] \writeline
    \item[$\square$] \writeline
\end{itemize}
\end{objectivebox}

\begin{vocabbox}
\termblanklong{$t$-test for slope}
\termblanklong{Null hypothesis for slope ($H_0$: $\beta = 0$)}
\termblanklong{Alternative hypothesis for slope ($H_a$)}
\termblanklong{$\beta_0$ (hypothesized slope value)}
\termblanklong{LINER conditions}
\end{vocabbox}

\begin{hookbox}
A researcher believes that elevation predicts average annual temperature for mountain
towns. She fits a regression line to data from $n = 18$ towns. Before running any
significance test, what three things does she need to set up?

\ans{(1) Name the procedure: $t$-test for slope. (2) State H$_0$ and H$_a$ in terms
of $\beta$. (3) Verify the LINER conditions.}
\end{hookbox}

\begin{notesbox}{Naming the Procedure and Stating Hypotheses (VAR-7.J/K)}
The appropriate procedure is a \ans{$t$-test} for slope.

\vspace{0.1in}
\textbf{Hypotheses} (always in terms of $\beta$, the \ans{population} slope):

\begin{itemize}
    \item H$_0$: $\beta = $ \ans{0} \quad (no linear relationship)
    \item H$_a$: $\beta$ \ans{$\ne$} 0 \; (two-sided) \quad OR \quad
          $\beta$ \ans{$>$} 0 \quad OR \quad $\beta$ \ans{$<$} 0
\end{itemize}

\vspace{0.05in}
\textit{Note:} Hypotheses are about $\beta$ (the \ans{population} slope),
not $b$ (the \ans{sample} slope).

\vspace{0.05in}
The hypothesized value $\beta_0 = 0$ because: \ans{we are testing whether there is
any linear relationship at all; $\beta = 0$ means no linear relationship exists.}
\end{notesbox}

\begin{notesbox}{LINER Conditions (VAR-7.L)}
\renewcommand{\arraystretch}{1.5}
\begin{tabularx}{\linewidth}{clX}
    \textbf{Letter} & \textbf{What It Means} & \textbf{How to Verify} \\
    \hline
    L & Linear relationship & Residual plot: \ans{random scatter, no curve} \\
    I & Independent observations & Random sample; $n \le$ \ans{10}$\%\cdot N$ \\
    N & Normal $y$-values & Residual \ans{histogram} or $n >$ \ans{30} \\
    E & Equal variance ($\sigma_y$ constant) & Residual plot: \ans{constant/equal spread} \\
    R & Random data source & \ans{random sample or randomized experiment} \\
\end{tabularx}

\vspace{0.1in}
\textbf{Practice Check:} ``A random sample of 30 students'' --- which conditions
can you verify from this statement alone? Which require the residual plot?

\ans{R (random sample) and I (10\% condition) can be checked from the statement;
L, N, and E require the residual plot.}
\end{notesbox}

\end{document}
